\subsection{Pipeline overview}
\label{sec:pipeline}

\begin{figure}[h]
\centering
\begin{tikzpicture}[
  every node/.style={font=\scriptsize, align=center},
  box/.style={rectangle, draw=black!70, line width=0.4pt,
              rounded corners=2pt, inner sep=2pt,
              minimum width=20mm, minimum height=8mm},
  inp/.style={box, fill=black!4},
  proc/.style={box, fill=black!8},
  arr/.style={-{Latex[length=1.6mm]}, line width=0.4pt},
  lbl/.style={midway, sloped, above, font=\tiny, inner sep=1pt}
]
% Left column: three input streams stacked vertically
\node[inp] (real)    at (0,  1.8) {Real\\examples};
\node[inp] (synth)   at (0,  0)   {Synthetic\\contrastive pairs};
\node[inp] (pool)    at (0, -1.8) {Unlabelled pool\\(enhanced only)};
% Bottom-row chain: leak audit + pseudo-labelled
\node[proc] (gate)   at (3.0, -1.8) {Leak audit\\\texttt{exp\_0081}};
\node[inp]  (pseudo) at (6.0, -1.8) {Pseudo-\\labelled};
% Central mixer
\node[proc] (mix)    at (6.0,  0)   {Mini-batch\\mix};
% Right pair: encoder, head
\node[proc] (enc)    at (9.0,  0)   {RoBERTa-large\\encoder};
\node[proc] (head)   at (12.0, 0)   {Per-task head\\(argmax /\\thresholds)};
% Arrows into mixer, with per-stream sw labels
\draw[arr] (real)   -- node[lbl] {\texttt{sw=1.0}} (mix);
\draw[arr] (synth)  -- node[midway, above, font=\tiny, inner sep=1pt] {\texttt{sw=0.5}} (mix);
\draw[arr] (pool)   -- (gate);
\draw[arr] (gate)   -- (pseudo);
\draw[arr] (pseudo) -- node[lbl] {\texttt{sw=0.7}} (mix);
% Right-side flow
\draw[arr] (mix)  -- (enc);
\draw[arr] (enc)  -- (head);
\end{tikzpicture}
\caption{Training and inference pipeline for the submitted RoBERTa-large systems.
Real, synthetic, and pseudo-labelled examples are mixed at per-stream loss weights
(1.0 / 0.5 / 0.7) into a single fine-tuning loop; the leak-audit gate sits between
pseudo-label generation and pseudo-label training. At inference time the trained
classifier head is decoded with per-class thresholds for ST3 and default argmax
for ST1/ST2.}
\label{fig:pipeline}
\end{figure}

The same pipeline produces all six submitted runs, with per-cell hyperparameters summarised in Appendix~\ref{app:hyper}. Real examples are concatenated from the published Reddit fields (\S\ref{sec:input}), augmented with a small batch of LLM-generated synthetic contrastive pairs (\S\ref{sec:synth}), and --- for ST2 and ST3 enhanced --- extended with pseudo-labelled unlabelled examples (\S\ref{sec:pseudo}) once a leakage audit gate has passed. The three input streams are interleaved in each training mini-batch under the per-stream loss weights given in the caption. The mixed corpus is fed into a RoBERTa-large encoder~\cite{liu2019roberta} with a task-specific classification head; loss is class-weighted cross-entropy for ST2 and ST3 (with class weights set to inverse-frequency on the dev partition) and unweighted cross-entropy for ST1. We use the HuggingFace \texttt{transformers} library throughout~\cite{wolf2020transformers}, with the AdamW optimiser, a linear warmup over the first 10\% of steps, and a peak learning rate of 1e-5. Inference for ST1 and ST2 uses default argmax over class logits; ST3 uses a per-class threshold sweep on the development partition (\S\ref{sec:thresholds}). The same training loop is used for all six (subtask, track) cells; only the input fields, the loss weights, the truncation length, and the presence or absence of pseudo-labelled examples vary across cells.

\subsection{Backbone selection}
\label{sec:backbone}

We compared RoBERTa-base, RoBERTa-large, DeBERTa-v3-base, DeBERTa-v3-large, and ModernBERT-large under matched data and training budgets. RoBERTa-large beat RoBERTa-base by 0.027 on ST1 base, 0.007 on ST1 enhanced, and 0.055 on ST2 base under matched recipes. % from exp_0013 exp_0004 exp_0022 exp_0016 exp_0014 exp_0005
The one cell where RoBERTa-base was marginally ahead --- ST2 enhanced (0.925 vs.\ 0.923) --- is inside seed-variance for these models and we treat it as a tie rather than evidence against the larger backbone. % from exp_0017 exp_0023
We therefore standardised on RoBERTa-large for the six submitted runs.

DeBERTa-v3-large did not survive the fine-tuning recipes we tried. Plain dev runs collapsed to near-chance F1 across all three subtasks (ST1 base 0.351, ST2 base 0.030, ST3 base 0.205 across two attempts). % from exp_0007 exp_0008 exp_0009 exp_0040
Mixed-precision training NaN-exploded across four subtask$\times$sw combinations; switching to bf16 with a smaller learning rate stopped the NaNs but produced near-zero F1. % from exp_0041 exp_0042 exp_0043 exp_0044 exp_0049 exp_0050 exp_0051
After four NaN-explosion incidents and seven collapse-or-fail runs, we removed DeBERTa-v3-large from the backbone shortlist.

ModernBERT-large~\cite{warner2024modernbert} ran cleanly but did not beat RoBERTa-large on ST3 enhanced, reaching 0.719 against RoBERTa-large's 0.735. % from exp_0086 exp_0073
A 1.6\,pp gap on a 4-way task is real but not enough to override the simpler, well-understood RoBERTa pipeline for a notebook-paper submission. For a corpus of roughly 473 ST3 examples, RoBERTa-large is the rational choice; we revisit the LLM alternatives in \S\ref{sec:llm-detour} and \S\ref{sec:disc-encoder}.

\subsection{Input construction}
\label{sec:input}

We concatenate the available text fields into a single input separated by the standard RoBERTa special tokens \texttt{<s>...</s></s>...</s>}. For the base track the order is \texttt{text\_raw}, \texttt{text\_raw\_parent}, \texttt{text\_raw\_title}, \texttt{text\_base}, \texttt{argument\_base}; for the enhanced track we append \texttt{text\_enhanced} and \texttt{argument\_enhanced}. Tokeniser truncation is set to \texttt{max\_len=384} for ST3 and \texttt{max\_len=256} for ST1 and ST2. % from exp_0069 exp_0073 exp_0033 exp_0035 exp_0032
The choice between 256 and 384 was driven by a token-length diagnostic showing that ST3 inputs benefited measurably from the longer window once \texttt{argument\_enhanced} was concatenated, while ST1/ST2 inputs fit comfortably inside 256 with only a small tail of truncations. % from exp_0061

\subsection{Synthetic data augmentation}
\label{sec:synth}

We generated synthetic contrastive pairs using a frontier LLM via API. For each target label, we prompted the model with three real in-domain examples and asked it to produce a paired contrast --- a minimally edited rewrite that flips the label while preserving topic, register, and Reddit-conversational style. Generations were filtered by a length sanity check and a string-overlap heuristic against the source pair to discourage trivial copies. The three batches are: \texttt{synth\_v002} (130 pairs for ST1/ST2, \texttt{exp\_0025}; 110 pairs for ST3, \texttt{exp\_0030}), % from exp_0025 exp_0030
\texttt{synth\_v003} (rare-class top-up for ST3, \texttt{exp\_0056}: 100 epistemic-external + 50 practical-external pairs), % from exp_0056
and the \textsf{pe}-targeted batch (\texttt{exp\_0103}: 100 practical-external + practical-anchor pairs).\footnote{The exact prompt template and the generated pairs are released alongside this notebook in the TIRA submission package.} % from exp_0103

Synth at \texttt{sw=0.5} added 0.014--0.019 F1 across the four ST1 and ST2 cells (Table~\ref{tab:ablation}). % from exp_0013 exp_0033 exp_0022 exp_0035 exp_0014 exp_0032 exp_0023 exp_0036
On ST3 enhanced, however, synth alone did not move the needle: the real-only baseline and the +synth configuration tied at 0.694, and the eventual gain on this cell only appeared once pseudo-labelling was added (\S\ref{sec:pseudo}). % from exp_0024 exp_0067

The largest single gain came on ST3 base, where \texttt{synth\_v002} at \texttt{max\_len=384} lifted F1 from 0.44 to 0.52, an improvement of +0.087. % from exp_0058 exp_0069
Synth is not a free lunch: adding the \texttt{synth\_v003} rare-class top-up on top of \texttt{synth\_v002} dropped ST3-base F1 by 0.045 relative to \texttt{synth\_v002} alone, and a focal-loss recipe~\cite{lin2017focal} lost 0.018 over the same baseline. % from exp_0068 exp_0069 exp_0053
We read this as confirmation of the general caveat documented by Li et al.~\cite{li2023synthdata}: synthetic data improves the average and harms the worst-case, and which side you land on depends on the recipe.

\subsection{Pseudo-labelling}
\label{sec:pseudo}

For ST2 and ST3 enhanced we additionally added pseudo-labelled examples drawn from the published unlabelled pool. Predictions from the best dev-F1 checkpoint at the matching subtask $\times$ track were filtered by confidence at threshold $\tau=0.9$ and added to the training mix at a loss weight \texttt{sw=0.7}.\footnote{The pseudo-label confidence threshold $\tau$ and the source checkpoint are released in the TIRA reproducibility package alongside the prompt template.} Following the modern transformer-era treatment by Mukherjee and Awadallah~\cite{mukherjee2020ust} and the original pseudo-label idea of Lee~\cite{lee2013pseudolabel}, we treated this step with caution: pseudo-labelling on enhanced-track inputs interacts directly with the leakage concern of \S\ref{sec:task-tracks}, so we ran a dedicated leak audit gate before letting pseudo examples into training; the audit returned no detectable leak under the test we ran and we proceeded. % from exp_0081
The recipe paid off most clearly on ST2 enhanced, where adding pseudo lifted F1 from 0.94 to 0.97, a gain of +0.033 over the synth-only baseline. % from exp_0036 exp_0072
On ST3 enhanced the same recipe at \texttt{max\_len=384} lifted F1 from the \texttt{synth\_v002}-only configuration to 0.74, a gain of +0.041 over the closest matching baseline. % from exp_0073
The pseudo recipe at \texttt{max\_len=256} underperformed its synth-only counterpart on ST1 enhanced (0.906 against 0.914, a loss of 0.008), % from exp_0071
which is one of the reasons we did not adopt pseudo for ST1 enhanced in the final submission.

\subsection{The LLM detour}
\label{sec:llm-detour}

Encoders did not look like the obvious choice at the start of the project. Following the published precedent of Pan et al.~\cite{pan2024zeroshotfallacy} --- whose zero-shot LLM results are reported on different fallacy corpora and are therefore not directly numerically comparable --- we ran a zero-shot Qwen2.5-3B baseline on ST3 enhanced, which reached F1 0.47 --- well below the RoBERTa baselines of \S\ref{sec:backbone}. % from exp_0091
A kNN few-shot variant retrieving over a \texttt{bge-large-en-v1.5} sentence embedding index (1876$\times$1024) recovered some ground at 0.53, but still trailed RoBERTa-large by a wide margin. % from exp_0138 exp_0143

We then trained CoT-LoRA recipes~\cite{hu2022lora,dettmers2023qlora,wei2022cot} on Qwen2.5-7B and Qwen2.5-32B. The 7B 5-fold CV on ST3 enhanced reached 0.85 $\pm$ 0.11; the 32B single-run reached 0.84, but the corresponding 5-fold CV was still running at the submission deadline. % from exp_0180 exp_0207
We also evaluated Qwen3-4B-Thinking~\cite{qwen3}: 5-fold CV reached 0.96 on ST2 enhanced, which is competitive with the submitted encoder, but on ST3 base the same model reached only 0.45, below the 0.52 we already had with RoBERTa-large + \texttt{synth\_v002}. % from exp_0217 exp_0218
Despite the apparent competitiveness of the larger CoT-LoRA configurations on individual cells, none of them delivered a validated improvement over the encoder on the cells we could test under cross-validation within our compute envelope. We treat the LLM track as a real route, not a dismissed baseline; the comparison is underpowered at our sample size.

\subsection{Threshold calibration for ST3 \textsf{pe}}
\label{sec:thresholds}

A first diagnostic on ST3 used an NLI-style basis probe whose per-class breakdown surfaced a clear weakness: \textsf{pe} F1 was 0.25, well below the other three classes. % from exp_0082
That number was what motivated the \textsf{pe}-targeted synth batch in \S\ref{sec:synth}. After training the final ST3-enhanced classifier, we ran a per-class threshold sweep on the dev partition; the best configuration --- $\{\text{ee}=0.1, \text{ei}=0.1, \text{pe}=0.05, \text{pi}=0.4\}$ --- added +1.0\,pp F1 over default argmax. % from exp_0145
A subsequent attempt to learn a basis-axis meta-learner from LLM signals gave no improvement over the encoder-only weighted average and was not used in the submission. % from exp_0146
